%% kapa: exam=1 round=23 year=2012 period=2 form=A subject=accounting subject_ko=회계학 seq=7 number=7
%% source: Q-Net 감정평가사 기출문제 (cst00309) artlSeq=5052580 file=raw/1008/5052580/2012년 제23회 감정평가사 제1차시험 기출문제(2교시 A형).pdf
%% conversion: kordoc 4.24.51 -> tools/gen1.py (auto); review=proofread-sonnet; flags=none
\begin{question}{7}
(주)강남은 사옥을 건설하기 위하여 20×2년 1월 1일에 (주)대한과 건설계약을 체결하였다. (주)강남의 사옥은 20×3년 6월 30일에 준공될 예정이고,
(주)강남은 사옥건설을 위해 다음과 같이 지출하였다. \par
\begin{center}
\small
\begin{tabular}{l c c c c}
일자 & \uline{20×2.1.1} & \uline{20×2.7.1} & \uline{20×2.10.1} & \uline{20×3.1.1} \\
금액 & ￦50,000 & ￦50,000 & ￦60,000 & ￦70,000 \\
\end{tabular}
\end{center}
(주)강남의 차입금은 다음과 같다. \par
\begin{center}
\small
\begin{tabular}{l l r l}
\uline{차입금} & \uline{차입일} & \uline{차입금액} & \uline{상환일} \\
K은행 & 20×2.1.1 & ￦50,000 & 20×3.6.30 \\
A은행 & 20×1.1.1 & ￦30,000 & 20×3.6.30 \\
B은행 & 20×1.1.1 & ￦50,000 & 20×4.12.31 \\
\end{tabular}

\medskip
\begin{tabular}{l c l}
\uline{차입금} & \uline{이자율} & \uline{이자지급조건} \\
K은행 & 12\% & 단리/매년말지급 \\
A은행 & 10\% & 단리/매년말지급 \\
B은행 & 12\% & 단리/매년말지급 \\
\end{tabular}
\end{center}
이들 차입금 중 K은행에서의 차입금은 (주)강남의 사옥건설을 위한 특정차입금이며, A은행 차입금과 B은행 차입금은 일반차입금이다. (주)강남의 건설 중인 사옥은 차입원가 자본화의 적격자산에 해당된다. 이에 대하여 (주)강남이 20×2년 자본화할 차입원가의 금액은? \par
\begin{choices}
  \item ￦4,500
  \item ￦6,000
  \item ￦9,000
  \item ￦10,500
  \item ￦13,500
\end{choices}
\end{question}
\answer{4}
